\chapter{Preliminaries}
\renewcommand{\currentchapterpath}{chapters/prelim/}

\begin{flushright}
\textit{“合抱之木，生于毫末；九层之台，起于累土。” ——《道德经》
\footnote{From the \textit{Dao De Jing}: ``A great tree grows from a tiny shoot; a nine-story terrace begins with a heap of earth.'' The quotation emphasizes that substantial achievements begin with small, steady steps.}}
\end{flushright}

This introduction chapter reviews some of the basic concepts and notions used
throughout the book. It aims to establish common vocabularies
and notations that are used in modern mathematical physics literature.
Though you, as a reader, may share little connection with me (the author),
it is hoped that after going through this chapter,
you and me are now connected in some way and we are on the same page.

A few notational conventions are used uniformly below.
The imaginary unit is $\mathrm{i}$ and the differential is $\mathrm{d}$,
both in upright roman type.
Vectors are written in bold, as in $\mathbf{A}$.
The Laplacian is $\nabla^2$; a finite increment is $\Delta x$.
A string or interval length is $\ell$.
The Fourier transform uses the symmetric convention
$F(\omega)=\mathcal{F}\{f\}(\omega)=(2\pi)^{-1/2}\int f(x)e^{-\mathrm{i}\omega x}\,\mathrm{d}x$.
Real logarithms are written $\ln$; the complex logarithm is $\log$.


 

\chapterinput{sets_and_logic}
\chapterinput{trigonometry}
\chapterinput{calculus}
\chapterinput{linear_algebra}
